% Generated draft: requires mathematical and source review.
\documentclass[11pt]{article}
\usepackage[margin=25mm]{geometry}
\usepackage{fontspec}
\setmainfont{Latin Modern Roman}
\usepackage{amsmath,amssymb,mathtools,unicode-math}
\setmathfont{Latin Modern Math}
\usepackage{enumitem,needspace,xurl,hyperref}
\setlength{\parindent}{0pt}
\setlength{\parskip}{5pt}
\setlength{\emergencystretch}{3em}
\newcommand{\N}{\mathbb{N}}
\newcommand{\Z}{\mathbb{Z}}
\newcommand{\Q}{\mathbb{Q}}
\newcommand{\R}{\mathbb{R}}
\newcommand{\C}{\mathbb{C}}
\newcommand{\F}{\mathbb{F}}
\newcommand{\A}{\mathbb{A}}
\newcommand{\PP}{\mathbb{P}}
\newcommand{\E}{\mathbb{E}}
\newcommand{\eps}{\varepsilon}
\newcommand{\dd}{\,\mathrm{d}}
\begin{document}
\section*{UnsolvedMath ID: 801112}
\textbf{Source problem code:} catalogue entry
\subsection{Definitions and mathematical statement}
Let \(\mathbb F_2=\langle a,b\rangle\) be the free group on two generators. Its group von Neumann algebra \(L(\mathbb F_2)\) is generated by the left regular representation on \(\ell^2(\mathbb F_2)\), with canonical trace \(\tau(x)=\langle x\delta_e,\delta_e\rangle\).

Under the usual interpretation in the continuous first-order language of tracial von Neumann algebras, a unital trace-preserving embedding \(\iota:A\to B\) is existential if, for every finite parameter tuple \(\bar a\) from \(A\) and every quantifier-free formula \(\varphi(\bar x,\bar y)\),
\[
\inf_{\bar y\in A}\varphi^A(\bar a,\bar y)
=
\inf_{\bar y\in B}\varphi^B(\iota(\bar a),\bar y),
\]
where the variables range over their specified operator-norm balls in both structures.

\textbf{Question 1.} Does there exist an existential embedding
\[
\iota:L(\mathbb F_2)\hookrightarrow\mathbb M^{\mathcal U}?
\]
Here \(\mathbb M^{\mathcal U}\) retains the source notation. The supplied record does not define the matrix algebras, their traces, the index set, or the ultrafilter used to construct this target. Thus this question requires the source's missing convention for \(\mathbb M^{\mathcal U}\); the interpretation of existentiality above is also subject to confirmation of the intended language.

\textbf{Question 2.} Is the source's random embedding of \(L(\mathbb F_2)\) into \(\mathbb M^{\mathcal U}\), described as arising from Haar-sampling, existential in this sense? Haar-sampling ordinarily refers to choosing unitaries according to normalized Haar probability measure on unitary groups. The record does not specify the matrix dimensions, joint distribution and independence assumptions, construction of the embedding, or intended probabilistic quantifier. In particular, it does not state whether existentiality is required almost surely. Consequently, the second question cannot be given a fully quantified probabilistic formulation from the supplied record alone.
\subsection{Short English statement}
The problem asks whether the free group factor \(L(\mathbb F_2)\) embeds existentially into the target denoted \(\mathbb M^{\mathcal U}\), and whether an embedding obtained by Haar-sampling is existential. Existentiality means that allowing witnesses in the larger algebra does not improve the infimum of any quantifier-free formula with parameters from the smaller algebra. The target construction and random sampling procedure require clarification.
\subsection{Sources}
\begin{itemize}
\item Local UnsolvedMath registry: permanent ID \(801112\), exact source problem code \texttt{catalogue entry}, title ``Ultraproduct II\(_1\) Factors --- Problem 13.'' The supplied record attributes the problem to Elayavalli, with proposed year 2025, and labels it open; this status is supplied metadata, not an independently verified assessment.
\item Supplied source list: Elayavalli, \emph{50 Open Problems: Ultraproduct II\(_1\) Factors} (2025), Problem 13, \url{https://arxiv.org/abs/2511.20377}. The record reports access on 2026-10-16 and extraction by \texttt{arxiv\_latex}; these are supplied provenance details, not a claim of independent consultation.
\item The supplied statement reports progress and evidence ``towards the above'' in the reference labeled \texttt{peterson2026ultraproductapproximationspropertyt}. No complete bibliographic entry or link for that reference is supplied, and the antecedent of ``the above'' is absent from the record.
\end{itemize}
\end{document}
